\chapter{Components}\label{ch:components}

This chapter describes every component type: what it models, how \app\ turns it into per
unit, which studies see it, and every one of its fields. The field tables are generated
from the program's own definitions. Each table gives the label shown on screen and,
below it, the field's name, which bindings, overrides and the data manager use. Fields
that depend on a mode appear in the property editor only when they apply, and the tables
say when.

Three fields are common to all components and are not repeated in the text below:
\gui{Name}, the label drawn on the diagram (the id stands in when it is empty), \gui{In
service}, cleared to make every study ignore the element, and the read-only \gui{id}. The
formulas of this chapter are collected, with the zero-sequence rules, in
\cref{app:modelling}.

\begin{note}
\app\ uses one base power $S_\text{base}$ for the whole network and one base voltage
$U_\text{base}$ per voltage zone (\cref{ch:bases}). ``Per unit'' below always means on
those bases. The base impedance at a bus is $Z_\text{base} = U_\text{base}^2 /
S_\text{base}$.
\end{note}

% ---------------------------------------------------------------------------
\componentsection{Bus}{bus}{A node of the network: a busbar, or a junction drawn as a dot.}

Buses are where everything connects. The \gui{Kind} only changes the drawing: a bar for a
\code{busbar}, a dot for a \code{junction}. Every bus gets its base voltage by propagation
from the base reference bus (\cref{ch:bases}). A value in \gui{Pinned U base} fixes the
base of that bus by hand instead. It also seeds the propagation, and a pinned base that
disagrees with a transformer ratio is reported as a conflict.

% Generated by tools/make_manual.py from psa.core.components.Bus.FIELDS. Do not edit.
\begin{fieldtable}
\fieldgroup{General}
\gui{Name}\newline\code{name} &  & -- & Label shown on the diagram. \\
\gui{In service}\newline\code{in\_\allowbreak{}service} &  & on & Out-of-service components are ignored by every study. \\
\fieldgroup{Parameters}
\gui{Kind}\newline\code{kind} &  & \code{busbar} & Junctions are drawn as a small dot instead of a bar. One of \code{busbar}, \code{junction}. \\
\gui{Pinned U base}\newline\code{u\_\allowbreak{}base\_\allowbreak{}kv} & kV & empty & Leave empty to derive the base by propagation from the reference bus through the transformer ratios. At least 0. \\
\end{fieldtable}


% ---------------------------------------------------------------------------
\componentsection{External grid}{source}{A voltage source behind a Thevenin impedance: the
network beyond the boundary.}

The external grid represents the system that feeds the network drawn. It is a voltage
$E = U\angle\theta$ behind the impedance $Z_\text{th}$. Give the impedance in one of
three ways (\gui{Impedance from}):

\begin{description}
  \item[\code{z\_ohm}] directly, in ohm, referred to the bus base. \code{0} makes an
    \emph{ideal} voltage source.
  \item[\code{ik}] from the three-phase fault current $I_k$ at the bus:
    $|Z_\text{th}| = U / (\sqrt3\, I_k)$.
  \item[\code{sk}] from the short-circuit power: $|Z_\text{th}| = U^2 / S_k$.
\end{description}

In the last two, $U = U_\text{pu}\cdot U_n$ with $U_n$ the \gui{U nominal} (the bus base
when it is empty), and the angle of the impedance follows from the power factor of the
fault: $R/|Z| = \cos\varphi$. The S2 grid, for instance, is 70~kV, 25~kA at
$\cos\varphi = 0.15$, which gives $Z_\text{th} = 0.2425 + j1.5983\,\Omega$.

How the studies see it:

\begin{itemize}
  \item \textbf{Linear solve, Thevenin, faults}: an ideal grid fixes its bus voltage. A
    grid behind an impedance becomes a Norton pair, $1/Z_\text{th}$ to earth plus the
    current $E/Z_\text{th}$. In the negative sequence it is $Z_\text{th}$ without a
    source voltage.
  \item \textbf{Admittance matrix}: depending on the study's \gui{Sources} option, the
    grid is an extra internal node behind $Z_\text{th}$, a shunt $1/Z_\text{th}$ to earth,
    or left out (\cref{sec:study-ybus}).
  \item \textbf{Load flow}: with \gui{Slack} checked, the grid is the slack bus at its own
    terminal, holding $U$ and $\theta$. Its impedance is short-circuit data and does not
    enter the load flow.
  \item \textbf{Zero sequence} (\gui{Z0 from}): \code{none}, an ungrounded system with no
    zero-sequence path; \code{z0\_ohm}, $Z_0$ given in ohm; \code{ik1}, from the single
    line-to-ground fault current, $Z_0 = 3U_\text{ph}/I_{k1}\angle\arccos(\cos\varphi_0) -
    2Z_1$; or \code{z0\_z1}, a multiple of $Z_1$.
\end{itemize}

% Generated by tools/make_manual.py from psa.core.components.ExternalGrid.FIELDS. Do not edit.
\begin{fieldtable}
\fieldgroup{General}
\gui{Name}\newline\code{name} &  & -- & Label shown on the diagram. \\
\gui{In service}\newline\code{in\_\allowbreak{}service} &  & on & Out-of-service components are ignored by every study. \\
\gui{Bus}\newline\code{bus} &  & -- & The bus this terminal is connected to: wire it on the canvas, or pick the bus here. \\
\fieldgroup{Parameters}
\gui{U}\newline\code{u\_\allowbreak{}pu} & pu & 1 & Source voltage; with Ik or Sk also the pre-fault voltage the fault level refers to. At least 0. \\
\gui{Angle}\newline\code{theta\_\allowbreak{}deg} & deg & 0 &  \\
\gui{Impedance from}\newline\code{z\_\allowbreak{}mode} &  & \code{z\_\allowbreak{}ohm} & z\_ohm: give Z directly. ik: three-phase fault current. sk: short-circuit power. One of \code{z\_\allowbreak{}ohm}, \code{ik}, \code{sk}. \\
\gui{Z Thevenin}\newline\code{z\_\allowbreak{}th\_\allowbreak{}ohm} & ohm & 0 & 0 means an ideal voltage source. Referred to the bus base. Shown when \gui{Impedance from} is \code{z\_\allowbreak{}ohm}. \\
\gui{Ik (three-phase)}\newline\code{ik\_\allowbreak{}ka} & kA & 0 & At least 0. Shown when \gui{Impedance from} is \code{ik}. \\
\gui{Sk (three-phase)}\newline\code{sk\_\allowbreak{}mva} & MVA & 0 & At least 0. Shown when \gui{Impedance from} is \code{sk}. \\
\gui{cos phi of the fault}\newline\code{cos\_\allowbreak{}phi} &  & 0.1 & Lagging power factor of the fault current: R/|Z|. From 0 to 1. Shown when \gui{Impedance from} is \code{ik} or \code{sk}. \\
\gui{U nominal}\newline\code{un\_\allowbreak{}kv} & kV & empty & Voltage the fault level is quoted at. Empty means the bus base. At least 0. Shown when \gui{Impedance from} is \code{ik} or \code{sk}. \\
\gui{Slack}\newline\code{is\_\allowbreak{}slack} &  & on & Fixes the voltage angle reference for the network. \\
\fieldgroup{Zero sequence}
\gui{Z0 from}\newline\code{z0\_\allowbreak{}mode} &  & \code{none} & none: no zero-sequence path (ungrounded). z0\_ohm: give Z0. ik1: single line-to-ground fault current. z0\_z1: Z0/Z1. One of \code{none}, \code{z0\_\allowbreak{}ohm}, \code{ik1}, \code{z0\_\allowbreak{}z1}. \\
\gui{Z0}\newline\code{z0\_\allowbreak{}ohm} & ohm & 0 & Shown when \gui{Z0 from} is \code{z0\_\allowbreak{}ohm}. \\
\gui{Ik1 (line-to-ground)}\newline\code{ik1\_\allowbreak{}ka} & kA & 0 & At least 0. Shown when \gui{Z0 from} is \code{ik1}. \\
\gui{cos phi of that fault}\newline\code{cos\_\allowbreak{}phi0} &  & 0.1 & From 0 to 1. Shown when \gui{Z0 from} is \code{ik1}. \\
\gui{Z0 / Z1}\newline\code{z0\_\allowbreak{}z1} &  & 1 & At least 0. Shown when \gui{Z0 from} is \code{z0\_\allowbreak{}z1}. \\
\end{fieldtable}


% ---------------------------------------------------------------------------
\componentsection{Generator}{generator}{A synchronous generator, for the load flow.}

The generator exists for the load flow. Its \gui{Control} sets what it holds:

\begin{description}
  \item[\code{pv}] its active power \gui{P} and its bus voltage \gui{U set}: the bus becomes
    a PU bus. \gui{Q min} and \gui{Q max} limit its reactive power when the load flow is
    asked to enforce limits; empty means no limit.
  \item[\code{slack}] the voltage \gui{U set} and its angle \gui{Angle}: the bus becomes
    the slack bus, which balances the network.
  \item[\code{pq}] a fixed \gui{P} and \gui{Q}: the bus stays a PQ bus.
\end{description}

The linear, Thevenin and fault studies do not see generators. They work with external
grids, so give a network an external grid when you want those studies too. Several
voltage-controlling devices on one bus (two generators, or a generator and an SVC) must
ask for the same voltage, otherwise the load flow refuses to run. They share the bus's
reactive power equally.

% Generated by tools/make_manual.py from psa.core.components.Generator.FIELDS. Do not edit.
\begin{fieldtable}
\fieldgroup{General}
\gui{Name}\newline\code{name} &  & -- & Label shown on the diagram. \\
\gui{In service}\newline\code{in\_\allowbreak{}service} &  & on & Out-of-service components are ignored by every study. \\
\gui{Bus}\newline\code{bus} &  & -- & The bus this terminal is connected to: wire it on the canvas, or pick the bus here. \\
\fieldgroup{Load flow}
\gui{Control}\newline\code{control} &  & \code{pv} & pv: holds P and the bus voltage. slack: holds the voltage and its angle and balances the network. pq: holds P and Q. One of \code{pv}, \code{slack}, \code{pq}. \\
\gui{P}\newline\code{p\_\allowbreak{}mw} & MW & 0 & Shown when \gui{Control} is \code{pv} or \code{pq}. \\
\gui{Q}\newline\code{q\_\allowbreak{}mvar} & MVAr & 0 & Shown when \gui{Control} is \code{pq}. \\
\gui{U set}\newline\code{u\_\allowbreak{}set\_\allowbreak{}pu} & pu & 1 & At least 0. Shown when \gui{Control} is \code{pv} or \code{slack}. \\
\gui{Angle}\newline\code{theta\_\allowbreak{}deg} & deg & 0 & Shown when \gui{Control} is \code{slack}. \\
\gui{Q min}\newline\code{q\_\allowbreak{}min\_\allowbreak{}mvar} & MVAr & empty & Empty: no limit. Shown when \gui{Control} is \code{pv}. \\
\gui{Q max}\newline\code{q\_\allowbreak{}max\_\allowbreak{}mvar} & MVAr & empty & Empty: no limit. Shown when \gui{Control} is \code{pv}. \\
\end{fieldtable}


% ---------------------------------------------------------------------------
\componentsection{Transformer}{transformer}{A two-winding transformer: a series impedance,
with the ratio absorbed by the base change.}

The transformer is its short-circuit impedance, $r$ and $x$ in percent on its own rating.
In per unit on the network bases,
\[
  z = \frac{r + jx}{100}\cdot\frac{S_\text{base}}{S_n}\cdot
      \left(\frac{U_n}{U_\text{base}}\right)^2,
\]
where $U_n$ is the rated voltage of the side chosen in \gui{Refer Z to} and
$U_\text{base}$ the base of the bus on that side. There is no magnetising branch. The
turns ratio does not appear because the bases on its two sides are in the rated ratio
(\cref{ch:bases}). When pinned bases make the two sides disagree, the per-unit table
reports the value the other side would give.

\begin{caution}
\gui{Tap ratio} is stored with the project but not modelled yet: an off-nominal tap gives
a warning, the transformer is treated as a plain series impedance, and the $\det F = 1$
check of the two-port study no longer applies.
\end{caution}

\paragraph{Zero sequence.} The \gui{Vector group} decides the zero-sequence path, with the
HV winding first: \code{YN} is a grounded star, \code{Y} an ungrounded star, \code{D} a
delta. With $Z_0$ from \gui{x0} and \gui{r0} (equal to $x$ and $r$ when empty), and $Z_n$
the neutral impedances, appearing as $3Z_n$:

\begin{itemize}
  \item \code{YNyn}: a series branch $Z_0 + 3Z_{n,\text{HV}} + 3Z_{n,\text{LV}}$;
  \item \code{YNd}: $Z_0 + 3Z_{n,\text{HV}}$ from the HV bus to earth, the LV side open;
  \item \code{Dyn}: $Z_0 + 3Z_{n,\text{LV}}$ from the LV bus to earth, the HV side open;
  \item every other group: open in the zero sequence.
\end{itemize}

Phase shifts (Yd11, Dyn5\dots) are not modelled, as in the course.

% Generated by tools/make_manual.py from psa.core.components.Transformer.FIELDS. Do not edit.
\begin{fieldtable}
\fieldgroup{General}
\gui{Name}\newline\code{name} &  & -- & Label shown on the diagram. \\
\gui{In service}\newline\code{in\_\allowbreak{}service} &  & on & Out-of-service components are ignored by every study. \\
\gui{HV bus}\newline\code{bus\_\allowbreak{}hv} &  & -- & The bus this terminal is connected to: wire it on the canvas, or pick the bus here. \\
\gui{LV bus}\newline\code{bus\_\allowbreak{}lv} &  & -- & The bus this terminal is connected to: wire it on the canvas, or pick the bus here. \\
\fieldgroup{Parameters}
\gui{S rated}\newline\code{sn\_\allowbreak{}mva} & MVA & 1 & At least 0. \\
\gui{U rated HV}\newline\code{u1n\_\allowbreak{}kv} & kV & 1 & At least 0. \\
\gui{U rated LV}\newline\code{u2n\_\allowbreak{}kv} & kV & 1 & At least 0. \\
\gui{x}\newline\code{x\_\allowbreak{}pct} & \% & 0 &  \\
\gui{r}\newline\code{r\_\allowbreak{}pct} & \% & 0 &  \\
\gui{Tap ratio}\newline\code{tap} &  & 1 & Off-nominal ratio. Must be 1.0 for the S1 studies. At least 0. \\
\gui{Refer Z to}\newline\code{referred\_\allowbreak{}side} &  & \code{lv} & Which side's rated voltage and base are used for the per-unit conversion. One of \code{lv}, \code{hv}. \\
\fieldgroup{Zero sequence}
\gui{Vector group}\newline\code{vector\_\allowbreak{}group} &  & \code{YNyn} & HV winding first: YN grounded star, Y star, D delta. Decides the zero-sequence path; phase shifts are not modelled. One of \code{YNyn}, \code{YNy}, \code{Yyn}, \code{Yy}, \code{YNd}, \code{Yd}, \code{Dyn}, \code{Dy}, \code{Dd}. \\
\gui{Zn HV}\newline\code{zn\_\allowbreak{}hv\_\allowbreak{}ohm} & ohm & 0 & Neutral-to-earth impedance of a grounded HV star (0 = solid). \\
\gui{Zn LV}\newline\code{zn\_\allowbreak{}lv\_\allowbreak{}ohm} & ohm & 0 & Neutral-to-earth impedance of a grounded LV star (0 = solid). \\
\gui{x0}\newline\code{x0\_\allowbreak{}pct} & \% & empty & Empty means the same as x. \\
\gui{r0}\newline\code{r0\_\allowbreak{}pct} & \% & empty & Empty means the same as r. \\
\end{fieldtable}


% ---------------------------------------------------------------------------
\componentsection{Line / Cable}{line}{A PI section, with $n$ identical lines in parallel.}

A line is a PI equivalent built from its per-kilometre data and its length $\ell$. With
$n$ lines in parallel,
\[
  z = \frac{(r + jx)\,\ell}{n\,Z_\text{base}}, \qquad
  \frac{y_\text{sh}}{2} = \frac{n\,(g + jb_c)\,\ell}{2}\,Z_\text{base}
  \quad\text{at each end.}
\]
\gui{b\_c} is the \emph{total} charging susceptance per kilometre of \emph{one} line. Half
of it goes to each end of the PI. A line must join two buses with the same base voltage.
Use a transformer to change voltage level.

\paragraph{Zero sequence.} \gui{Z0 data} must be set before any zero-sequence study runs:
\code{factors} multiplies $r$, $x$ (and $g$, $b_c$ by \gui{b0 / b}), and \code{absolute}
gives $r_0$, $x_0$, $b_0$ per kilometre. While it is \code{unset}, the Thevenin study
leaves the zero sequence out with a warning, and the studies that need it (single
line-to-ground and double line-to-ground faults, unbalanced loads, the zero-sequence
admittance matrix) refuse to run and name the line.

% Generated by tools/make_manual.py from psa.core.components.Line.FIELDS. Do not edit.
\begin{fieldtable}
\fieldgroup{General}
\gui{Name}\newline\code{name} &  & -- & Label shown on the diagram. \\
\gui{In service}\newline\code{in\_\allowbreak{}service} &  & on & Out-of-service components are ignored by every study. \\
\gui{From bus}\newline\code{bus\_\allowbreak{}from} &  & -- & The bus this terminal is connected to: wire it on the canvas, or pick the bus here. \\
\gui{To bus}\newline\code{bus\_\allowbreak{}to} &  & -- & The bus this terminal is connected to: wire it on the canvas, or pick the bus here. \\
\fieldgroup{Parameters}
\gui{r}\newline\code{r\_\allowbreak{}ohm\_\allowbreak{}per\_\allowbreak{}km} & ohm/km & 0 &  \\
\gui{x}\newline\code{x\_\allowbreak{}ohm\_\allowbreak{}per\_\allowbreak{}km} & ohm/km & 0 &  \\
\gui{b\_c}\newline\code{b\_\allowbreak{}s\_\allowbreak{}per\_\allowbreak{}km} & S/km & 0 & Total charging susceptance per km of ONE line. \\
\gui{g}\newline\code{g\_\allowbreak{}s\_\allowbreak{}per\_\allowbreak{}km} & S/km & 0 &  \\
\gui{Length}\newline\code{length\_\allowbreak{}km} & km & 1 & At least 0. \\
\gui{Parallel lines}\newline\code{n\_\allowbreak{}parallel} &  & \code{1} & At least 1. \\
\fieldgroup{Zero sequence}
\gui{Z0 data}\newline\code{z0\_\allowbreak{}mode} &  & \code{unset} & unset: zero-sequence studies refuse to run. factors: multiples of r, x, b. absolute: per-km values. One of \code{unset}, \code{factors}, \code{absolute}. \\
\gui{r0 / r}\newline\code{r0\_\allowbreak{}factor} &  & 3 & At least 0. Shown when \gui{Z0 data} is \code{factors}. \\
\gui{x0 / x}\newline\code{x0\_\allowbreak{}factor} &  & 3 & At least 0. Shown when \gui{Z0 data} is \code{factors}. \\
\gui{b0 / b}\newline\code{b0\_\allowbreak{}factor} &  & 0.5 & At least 0. Shown when \gui{Z0 data} is \code{factors}. \\
\gui{r0}\newline\code{r0\_\allowbreak{}ohm\_\allowbreak{}per\_\allowbreak{}km} & ohm/km & 0 & At least 0. Shown when \gui{Z0 data} is \code{absolute}. \\
\gui{x0}\newline\code{x0\_\allowbreak{}ohm\_\allowbreak{}per\_\allowbreak{}km} & ohm/km & 0 & Shown when \gui{Z0 data} is \code{absolute}. \\
\gui{b0}\newline\code{b0\_\allowbreak{}s\_\allowbreak{}per\_\allowbreak{}km} & S/km & 0 & Shown when \gui{Z0 data} is \code{absolute}. \\
\end{fieldtable}


% ---------------------------------------------------------------------------
\componentsection{Load}{load}{A load given by its power at a nominal voltage.}

Give the load's power as \gui{P} with \gui{cos phi} and \gui{Power factor} (\code{lag}
draws inductive reactive power, \code{lead} capacitive), or as \gui{P} and \gui{Q}
directly (\gui{Given as} = \code{p\_q}, positive $Q$ inductive). Its \gui{Model} decides
how the studies see it:

\begin{description}
  \item[\code{constant\_impedance}] the load is the admittance that draws its power at
    \gui{U nominal} (the bus base when empty):
    $y = \dfrac{\conj{S}}{S_\text{base}}\left(\dfrac{U_\text{base}}{U_n}\right)^2$.
    A load quoted at 12~kV on a bus whose base is 12.6~kV therefore draws more than its
    nominal power at 1~pu, as a real impedance load would.
  \item[\code{constant\_power}] in the load flow, the load is a PQ injection that draws
    exactly its power whatever the voltage. The linear, Thevenin and fault studies cannot
    represent constant power: they keep the constant impedance and say so in a warning.
\end{description}

\paragraph{Zero sequence and unbalance.} \gui{Connection} \code{Y0} is a grounded star,
whose zero-sequence admittance is $1/(1/y + 3Z_n)$. \code{Y} and \code{D} are open in the
zero sequence. The unbalanced load study (\cref{sec:study-unbalanced}) uses the connection
to build the load's phase admittance matrix.

% Generated by tools/make_manual.py from psa.core.components.Load.FIELDS. Do not edit.
\begin{fieldtable}
\fieldgroup{General}
\gui{Name}\newline\code{name} &  & -- & Label shown on the diagram. \\
\gui{In service}\newline\code{in\_\allowbreak{}service} &  & on & Out-of-service components are ignored by every study. \\
\gui{Bus}\newline\code{bus} &  & -- & The bus this terminal is connected to: wire it on the canvas, or pick the bus here. \\
\fieldgroup{Parameters}
\gui{P}\newline\code{p\_\allowbreak{}mw} & MW & 0 &  \\
\gui{Given as}\newline\code{spec} &  & \code{p\_\allowbreak{}cosphi} & P and cos phi, or P and Q. One of \code{p\_\allowbreak{}cosphi}, \code{p\_\allowbreak{}q}. \\
\gui{cos phi}\newline\code{cos\_\allowbreak{}phi} &  & 1 & From 0 to 1. Shown when \gui{Given as} is \code{p\_\allowbreak{}cosphi}. \\
\gui{Power factor}\newline\code{pf\_\allowbreak{}sign} &  & \code{lag} & One of \code{lag}, \code{lead}. Shown when \gui{Given as} is \code{p\_\allowbreak{}cosphi}. \\
\gui{Q}\newline\code{q\_\allowbreak{}mvar} & MVAr & 0 & Positive is inductive (the load absorbs reactive power). Shown when \gui{Given as} is \code{p\_\allowbreak{}q}. \\
\gui{U nominal}\newline\code{u\_\allowbreak{}nominal\_\allowbreak{}kv} & kV & empty & Voltage at which the load consumes P. Empty means the bus base. At least 0. \\
\gui{Model}\newline\code{model} &  & \code{constant\_\allowbreak{}impedance} & The load flow honours constant\_power (a PQ injection); the linear, Thevenin and fault studies treat every load as a constant impedance at its nominal voltage. One of \code{constant\_\allowbreak{}impedance}, \code{constant\_\allowbreak{}power}. \\
\fieldgroup{Zero sequence}
\gui{Connection}\newline\code{connection} &  & \code{Y0} & Y0: grounded star (a zero-sequence path). Y and D: none. One of \code{Y0}, \code{Y}, \code{D}. \\
\gui{Zn}\newline\code{zn\_\allowbreak{}ohm} & ohm & 0 & Neutral-to-earth impedance (0 = solidly grounded). Shown when \gui{Connection} is \code{Y0}. \\
\end{fieldtable}


% ---------------------------------------------------------------------------
\componentsection{Shunt}{shunt}{A shunt capacitor ($Q>0$) or reactor ($Q<0$).}

The shunt is rated by the reactive power $Q$ it delivers at \gui{U nominal}. Positive $Q$
is capacitive, the same sign as line charging. Use a negative value for a reactor; the
symbol then changes to a coil. As an admittance,
$y = j\,\dfrac{Q}{S_\text{base}}\left(\dfrac{U_\text{base}}{U_n}\right)^2$.

In the load flow, \gui{Model} chooses between \code{constant\_b}, the admittance above,
which delivers $B\,U^2$, and \code{constant\_q}, an injection of $Q$ whatever the voltage.
The two readings are the base case and the variant of Assignment LF. Every other study
uses the admittance. \gui{Connection} works as for loads.

% Generated by tools/make_manual.py from psa.core.components.Shunt.FIELDS. Do not edit.
\begin{fieldtable}
\fieldgroup{General}
\gui{Name}\newline\code{name} &  & -- & Label shown on the diagram. \\
\gui{In service}\newline\code{in\_\allowbreak{}service} &  & on & Out-of-service components are ignored by every study. \\
\gui{Bus}\newline\code{bus} &  & -- & The bus this terminal is connected to: wire it on the canvas, or pick the bus here. \\
\fieldgroup{Parameters}
\gui{Q}\newline\code{q\_\allowbreak{}mvar} & MVAr & 0 & Positive is capacitive, the same sign as line charging. Use a negative value for a reactor. \\
\gui{U nominal}\newline\code{u\_\allowbreak{}nominal\_\allowbreak{}kv} & kV & empty & At least 0. \\
\fieldgroup{Load flow}
\gui{Model}\newline\code{model} &  & \code{constant\_\allowbreak{}b} & constant\_b: Q is rated at U nominal and scales with U\textasciicircum{}2. constant\_q: the load flow injects Q whatever the voltage. One of \code{constant\_\allowbreak{}b}, \code{constant\_\allowbreak{}q}. \\
\fieldgroup{Zero sequence}
\gui{Connection}\newline\code{connection} &  & \code{Y0} & Y0: grounded star (a zero-sequence path). Y and D: none. One of \code{Y0}, \code{Y}, \code{D}. \\
\end{fieldtable}


% ---------------------------------------------------------------------------
\componentsection{SVC}{svc}{A voltage-controlled shunt: a static var compensator or a
synchronous condenser.}

In the load flow, the SVC holds its bus at \gui{U set} with whatever reactive power that
takes, so its bus is a PU bus with $P = 0$. The results report that $Q$, and the
susceptance it amounts to, $B = Q/U^2$, in per unit and in siemens. This is how
Assignment LF sizes the controllable shunt at bus~3. \gui{Q min} and \gui{Q max} limit it
when the load flow enforces limits. The other studies do not see the SVC.

% Generated by tools/make_manual.py from psa.core.components.SVC.FIELDS. Do not edit.
\begin{fieldtable}
\fieldgroup{General}
\gui{Name}\newline\code{name} &  & -- & Label shown on the diagram. \\
\gui{In service}\newline\code{in\_\allowbreak{}service} &  & on & Out-of-service components are ignored by every study. \\
\gui{Bus}\newline\code{bus} &  & -- & The bus this terminal is connected to: wire it on the canvas, or pick the bus here. \\
\fieldgroup{Load flow}
\gui{U set}\newline\code{u\_\allowbreak{}set\_\allowbreak{}pu} & pu & 1 & At least 0. \\
\gui{Q min}\newline\code{q\_\allowbreak{}min\_\allowbreak{}mvar} & MVAr & empty & Empty: no limit (negative is inductive). \\
\gui{Q max}\newline\code{q\_\allowbreak{}max\_\allowbreak{}mvar} & MVAr & empty & Empty: no limit. \\
\end{fieldtable}


% ---------------------------------------------------------------------------
\componentsection{Impedance}{impedance}{A raw $R + jX$, in series between two buses or to
earth.}

A plain impedance, for whatever the other components do not cover: a series reactor, a
fault path, an equivalent. Give $R$ and $X$ in ohm, referred to the base of the
\gui{From bus}, or check \gui{Values in pu}. With \gui{To bus} empty, the impedance goes
from its bus to earth. It is the same in every sequence.

% Generated by tools/make_manual.py from psa.core.components.Impedance.FIELDS. Do not edit.
\begin{fieldtable}
\fieldgroup{General}
\gui{Name}\newline\code{name} &  & -- & Label shown on the diagram. \\
\gui{In service}\newline\code{in\_\allowbreak{}service} &  & on & Out-of-service components are ignored by every study. \\
\gui{From bus}\newline\code{bus\_\allowbreak{}from} &  & -- & The bus this terminal is connected to: wire it on the canvas, or pick the bus here. \\
\gui{To bus}\newline\code{bus\_\allowbreak{}to} &  & empty & Leave empty for a shunt to ground. \\
\fieldgroup{Parameters}
\gui{R}\newline\code{r} & ohm / pu & 0 &  \\
\gui{X}\newline\code{x} & ohm / pu & 0 &  \\
\gui{Values in pu}\newline\code{in\_\allowbreak{}pu} &  & off &  \\
\end{fieldtable}

